\chapter[Lezione 12 --- 22 ottobre 2026]{Lezione 12: teoremi del confronto e teorema delle tre funzioni}
\label{chap:lezione-12}

\section{Primo teorema del confronto}
\label{sec:lez12-confronto-primo}

\begin{theorem}[Primo teorema del confronto]
\label{thm:lez12-confronto-primo}
Siano \(f,g\colon A\longrightarrow\R\), con
\(A\subseteq\R\) non vuoto, e sia
\(x_0\in\overline{\R}\) un punto di accumulazione per \(A\).
Valgono le seguenti affermazioni.
\begin{enumerate}
  \item Se
  \[
    \lim_{x\to x_0}f(x)=\ell_1\in\R,
    \qquad
    \lim_{x\to x_0}g(x)=\ell_2\in\R,
    \qquad \ell_1<\ell_2,
  \]
  allora esiste un intorno \(U_{x_0}\) di \(x_0\) tale che
  \[
    f(x)<g(x)
    \qquad\forall x\in(U_{x_0}\cap A)\setminus\{x_0\}.
  \]
  \item Se
  \[
    \lim_{x\to x_0}f(x)=-\infty,
    \qquad
    \lim_{x\to x_0}g(x)=\ell\in\R\cup\{+\infty\},
  \]
  allora esiste un intorno \(U_{x_0}\) di \(x_0\) tale che
  \[
    f(x)<g(x)
    \qquad\forall x\in(U_{x_0}\cap A)\setminus\{x_0\}.
  \]
  \item Se
  \[
    \lim_{x\to x_0}f(x)=\ell\in\R,
    \qquad
    \lim_{x\to x_0}g(x)=+\infty,
  \]
  allora esiste un intorno \(U_{x_0}\) di \(x_0\) tale che
  \[
    f(x)<g(x)
    \qquad\forall x\in(U_{x_0}\cap A)\setminus\{x_0\}.
  \]
\end{enumerate}
\end{theorem}

\begin{proof}[Dimostrazione del primo punto]
Poiché \(\ell_1<\ell_2\), si ponga
\[
  \varepsilon:=\frac{\ell_2-\ell_1}{2}>0.
\]
Per la \cref{def:lez08-limite}, esistono intorni
\(U'_{x_0}\) e \(U''_{x_0}\) di \(x_0\) tali che
\[
  |f(x)-\ell_1|<\varepsilon
  \qquad\forall x\in(U'_{x_0}\cap A)\setminus\{x_0\},
\]
\[
  |g(x)-\ell_2|<\varepsilon
  \qquad\forall x\in(U''_{x_0}\cap A)\setminus\{x_0\}.
\]
Per la \cref{sec:lez11-intersezione-intorni},
\[
  U_{x_0}:=U'_{x_0}\cap U''_{x_0}
\]
è ancora un intorno di \(x_0\). Per ogni
\(x\in(U_{x_0}\cap A)\setminus\{x_0\}\), si ottiene
\[
  f(x)<\ell_1+\varepsilon
  =\frac{\ell_1+\ell_2}{2}
  =\ell_2-\varepsilon<g(x).
\]
\end{proof}

\clearpage

\begin{exercise}[Casi con limiti infiniti]
\label{ex:lez12-confronto-primo-infiniti}
Dimostrare il secondo e il terzo punto del
\cref{thm:lez12-confronto-primo} mediante la
\cref{def:lez08-limite}.
\end{exercise}

\begin{proof}[Soluzione dell'esercizio]
\emph{Secondo punto.}
Se \(\ell\in\R\), dalla \cref{def:lez08-limite},
scegliendo \(\varepsilon=1\), esiste un intorno
\(U'_{x_0}\) nel quale \(g(x)>\ell-1\).
Poiché \(f(x)\to-\infty\), esiste un intorno
\(U''_{x_0}\) nel quale \(f(x)<\ell-1\).
Per ogni \(x\in(U'_{x_0}\cap U''_{x_0}\cap A)
\setminus\{x_0\}\), segue
\[
  f(x)<\ell-1<g(x).
\]
Se \(\ell=+\infty\), esistono intorni
\(U'_{x_0}\) e \(U''_{x_0}\) nei quali valgono
rispettivamente \(f(x)<0\) e \(g(x)>0\).
Per ogni \(x\in(U'_{x_0}\cap U''_{x_0}\cap A)\setminus\{x_0\}\),
si ha
\[
  f(x)<0<g(x).
\]

\emph{Terzo punto.}
Poiché \(f(x)\to\ell\in\R\), scegliendo
\(\varepsilon=1\) esiste un intorno \(U'_{x_0}\)
nel quale \(f(x)<\ell+1\).
Poiché \(g(x)\to+\infty\), esiste un intorno
\(U''_{x_0}\) nel quale \(g(x)>\ell+1\).
Per ogni \(x\in(U'_{x_0}\cap U''_{x_0}\cap A)
\setminus\{x_0\}\), segue
\[
  f(x)<\ell+1<g(x).
\]
In ciascun caso, l'intersezione dei due intorni è un
intorno di \(x_0\) per la \cref{sec:lez11-intersezione-intorni}.
\end{proof}

\begin{example}[Limiti uguali e ordinamento variabile]
\label{ex:lez12-limiti-uguali-ordine}
Si considerino le funzioni su \(\R\)
\[
  f(x)=x^3,
  \qquad g(x)=-x^3.
\]
Entrambe hanno limite \(0\) per \(x\to0\):
per ogni \(\varepsilon>0\), scegliendo
\(\delta=\varepsilon^{1/3}\), da \(0<|x|<\delta\)
segue \(|x^3|<\varepsilon\).
Tuttavia,
\[
  f(x)<g(x)\quad\text{se }x<0,
  \qquad
  f(x)>g(x)\quad\text{se }x>0.
\]
Ogni intorno bucato di \(0\) contiene punti nei quali
si verifica ciascuna delle due disuguaglianze.
\end{example}

\begin{figure}[htbp]
  \centering
  \begin{tikzpicture}[unipd/diagram,x=2.25cm,y=1.28cm]
  \draw[->,unipd-muted] (-1.55,0) -- (1.6,0) node[right] {\(x\)};
  \draw[->,unipd-muted] (0,-2.75) -- (0,2.8) node[above] {\(y\)};
  \draw[unipd-primary,line width=1.05pt]
    plot[domain=-1.33:1.33,samples=100,smooth]
      (\x,{\x*\x*\x});
  \draw[unipd-accent,line width=1.05pt]
    plot[domain=-1.33:1.33,samples=100,smooth]
      (\x,{-\x*\x*\x});
  \fill[unipd-muted] (0,0) circle (1.7pt);
  \node[anchor=west,text=unipd-primary] at (1.03,1.78) {\(f(x)=x^3\)};
  \node[anchor=west,text=unipd-accent] at (1.03,-1.78) {\(g(x)=-x^3\)};
  \node[below left] at (0,0) {\(0\)};
\end{tikzpicture}

  \caption{Le funzioni \(x^3\) e \(-x^3\) hanno lo
  stesso limite per \(x\to0\), ma il loro ordine si
  inverte attraversando l'origine.}
  \label{fig:lez12-limiti-uguali-ordine}
\end{figure}

\section{Secondo teorema del confronto}
\label{sec:lez12-confronto-secondo}

\begin{theorem}[Secondo teorema del confronto]
\label{thm:lez12-confronto-secondo}
Siano \(f,g\colon A\longrightarrow\R\), con
\(A\subseteq\R\) non vuoto, e sia
\(x_0\in\overline{\R}\) un punto di accumulazione per \(A\).
Si supponga che esista un intorno \(U_{x_0}\) di
\(x_0\) tale che
\[
  f(x)\le g(x)
  \qquad\forall x\in(U_{x_0}\cap A)\setminus\{x_0\}.
\]
Allora:
\begin{enumerate}
  \item se
  \[
    \lim_{x\to x_0}f(x)=\ell_1\in\R,
    \qquad
    \lim_{x\to x_0}g(x)=\ell_2\in\R,
  \]
  si ha \(\ell_1\le\ell_2\);
  \item se \(\displaystyle\lim_{x\to x_0}g(x)=-\infty\),
  allora \(\displaystyle\lim_{x\to x_0}f(x)=-\infty\);
  \item se \(\displaystyle\lim_{x\to x_0}f(x)=+\infty\),
  allora \(\displaystyle\lim_{x\to x_0}g(x)=+\infty\).
\end{enumerate}
\end{theorem}

La dimostrazione del \cref{thm:lez12-confronto-secondo}
non viene sviluppata in questa sede.

La disuguaglianza stretta \(f(x)<g(x)\)
non garantisce \(\ell_1<\ell_2\).

\begin{example}[Disuguaglianza stretta e limiti uguali]
\label{ex:lez12-confronto-stretto-limiti-uguali}
Si considerino le funzioni \(f,g\colon\R\longrightarrow\R\)
definite da
\[
  f(x):=0,
  \qquad
  g(x):=
  \begin{cases}
    \dfrac1x, & x>0,\\[4pt]
    0, & x=0,\\[4pt]
    -\dfrac1x, & x<0.
  \end{cases}
\]
Si ha \(f(x)<g(x)\) per ogni \(x\ne0\). Inoltre,
per ogni \(\varepsilon>0\), scegliendo \(b=1/\varepsilon\),
per \(x>b\) si ottiene \(0<g(x)=1/x<\varepsilon\).
Pertanto
\[
  \lim_{x\to+\infty}f(x)
  =\lim_{x\to+\infty}g(x)=0.
\]
\end{example}

\begin{figure}[htbp]
  \centering
  \begin{tikzpicture}[unipd/diagram,x=.86cm,y=1.05cm]
  \draw[->,unipd-muted] (-5.1,0) -- (5.35,0) node[right] {\(x\)};
  \draw[->,unipd-muted] (0,-.3) -- (0,3.45) node[above] {\(y\)};
  \draw[unipd-accent,line width=1.1pt] (-4.95,0) -- (5.1,0);
  \draw[unipd-primary,line width=1pt]
    plot[domain=-4.9:-.36,samples=120,smooth]
      (\x,{-1/\x});
  \draw[unipd-primary,line width=1pt]
    plot[domain=.36:4.9,samples=120,smooth]
      (\x,{1/\x});
  \fill[unipd-primary] (0,0) circle (1.8pt);
  \node[text=unipd-primary,anchor=west] at (1.25,1.45)
    {\(g(x)=1/|x|\quad(x\ne0)\)};
  \node[text=unipd-accent,anchor=west] at (2.45,-.23)
    {\(f(x)=0\)};
  \node[below left] at (0,0) {\(0\)};
\end{tikzpicture}

  \caption{La funzione \(g(x)>f(x)=0\) per ogni
  \(x\ne0\), ma entrambe hanno limite \(0\) per
  \(x\to+\infty\).}
  \label{fig:lez12-stretto-limiti-uguali}
\end{figure}

\clearpage

\begin{example}[Un limite infinito mediante confronto]
\label{ex:lez12-limite-x-seno}
Si consideri la funzione
\[
  g\colon\R\longrightarrow\R,
  \qquad g(x)=x+\sin x.
\]
Poiché \(\sin x\ge-1\) per ogni \(x\in\R\), si ha
\[
  x-1\le x+\sin x
  \qquad\forall x\in\R.
\]
Per \(f(x):=x-1\) si ha
\[
  \lim_{x\to+\infty}f(x)=+\infty,
\]
poiché, per ogni \(K\in\R\), la condizione \(x>K+1\)
implica \(f(x)>K\). Per il terzo punto del
\cref{thm:lez12-confronto-secondo}, segue
\[
  \lim_{x\to+\infty}(x+\sin x)=+\infty.
\]
\end{example}

\begin{figure}[htbp]
  \centering
  \begin{tikzpicture}[unipd/diagram,x=.92cm,y=.62cm]
  \draw[->,unipd-muted] (-.35,0) -- (7.2,0) node[right] {\(x\)};
  \draw[->,unipd-muted] (0,-1.25) -- (0,8.25) node[above] {\(y\)};
  \draw[unipd-muted,densely dashed,line width=.75pt]
    plot[domain=0:6.25,samples=2] (\x,{\x+1});
  \draw[unipd-accent,line width=1pt]
    plot[domain=0:6.25,samples=2] (\x,{\x-1});
  \draw[unipd-primary,line width=1.15pt]
    plot[domain=0:6.25,samples=180,smooth]
      (\x,{\x+sin(deg(\x))});
  \node[font=\sffamily\footnotesize,text=unipd-muted,anchor=west]
    at (6.3,7.25) {\(x+1\)};
  \node[font=\sffamily\footnotesize,text=unipd-primary,anchor=west]
    at (6.3,6.22) {\(x+\sin x\)};
  \node[font=\sffamily\footnotesize,text=unipd-accent,anchor=west]
    at (6.3,5.25) {\(x-1\)};
  \node[below left] at (0,0) {\(0\)};
\end{tikzpicture}

  \caption{La funzione \(x+\sin x\) è compresa tra
  \(x-1\) e \(x+1\). In particolare, il minorante
  \(x-1\) tende a \(+\infty\).}
  \label{fig:lez12-confronto-seno}
\end{figure}

\section{Teorema delle tre funzioni}
\label{sec:lez12-tre-funzioni}

\begin{theorem}[Teorema delle tre funzioni]
\label{thm:lez12-tre-funzioni}
Siano \(f,g,h\colon A\longrightarrow\R\), con
\(A\subseteq\R\) non vuoto, e sia
\(x_0\in\overline{\R}\) un punto di accumulazione per \(A\).
Si supponga che
\[
  \lim_{x\to x_0}f(x)=\ell\in\R,
  \qquad
  \lim_{x\to x_0}g(x)=\ell,
\]
e che esista un intorno \(U_{x_0}\) di \(x_0\) tale che
\[
  f(x)\le h(x)\le g(x)
  \qquad\forall x\in(U_{x_0}\cap A)\setminus\{x_0\}.
\]
Allora
\[
  \lim_{x\to x_0}h(x)=\ell.
\]
\end{theorem}

\begin{proof}
Sia \(\varepsilon>0\). Per la \cref{def:lez08-limite},
esistono intorni \(U'_{x_0}\) e \(U''_{x_0}\) di
\(x_0\) tali che
\[
  |f(x)-\ell|<\varepsilon
  \qquad\forall x\in(U'_{x_0}\cap A)\setminus\{x_0\},
\]
\[
  |g(x)-\ell|<\varepsilon
  \qquad\forall x\in(U''_{x_0}\cap A)\setminus\{x_0\}.
\]
Per la \cref{sec:lez11-intersezione-intorni},
\[
  W_{x_0}:=U'_{x_0}\cap U''_{x_0}\cap U_{x_0}
\]
è ancora un intorno di \(x_0\). Per ogni
\(x\in(W_{x_0}\cap A)\setminus\{x_0\}\) valgono
simultaneamente le disuguaglianze
\[
  \ell-\varepsilon<f(x)\le h(x)
  \le g(x)<\ell+\varepsilon.
\]
Quindi
\[
  |h(x)-\ell|<\varepsilon
  \qquad\forall x\in(W_{x_0}\cap A)\setminus\{x_0\},
\]
e, per la \cref{def:lez08-limite}, si conclude
\(\lim_{x\to x_0}h(x)=\ell\).
\end{proof}

\begin{exercise}[Limite di \(x\sin(1/x)\)]
\label{ex:lez12-limite-seno-inverso}
Verificare mediante il \cref{thm:lez12-tre-funzioni} che
\[
  \lim_{x\to0}\left|x\sin\frac1x\right|=0,
  \qquad
  \lim_{x\to0}x\sin\frac1x=0.
\]
\end{exercise}

\begin{figure}[htbp]
  \centering
  \begin{tikzpicture}[unipd/diagram,x=4.5cm,y=2.9cm]
  \draw[->,unipd-muted] (-1.12,0) -- (1.15,0) node[right] {\(x\)};
  \draw[->,unipd-muted] (0,-1.08) -- (0,1.09) node[above] {\(y\)};
  \draw[unipd-accent,densely dashed,line width=.85pt]
    (-.95,.95) -- (0,0) -- (.95,.95);
  \draw[unipd-accent,densely dashed,line width=.85pt]
    (-.95,-.95) -- (0,0) -- (.95,-.95);
  \draw[unipd-primary,line width=1pt]
    plot[domain=-.95:-.032,samples=350,smooth]
      (\x,{\x*sin(deg(1/\x))});
  \draw[unipd-primary,line width=1pt]
    plot[domain=.032:.95,samples=350,smooth]
      (\x,{\x*sin(deg(1/\x))});
  \fill[white] (0,0) circle (2.2pt);
  \draw[unipd-primary,line width=.85pt] (0,0) circle (2.2pt);
  \node[font=\sffamily\footnotesize,text=unipd-accent,anchor=west]
    at (.61,.86) {\(|x|\)};
  \node[font=\sffamily\footnotesize,text=unipd-accent,anchor=west]
    at (.62,-.88) {\(-|x|\)};
\end{tikzpicture}

  \caption{La funzione \(x\sin(1/x)\), non definita in
  \(x=0\), rimane compresa tra \(-|x|\) e \(|x|\),
  entrambe convergenti a \(0\) per \(x\to0\).}
  \label{fig:lez12-seno-inverso}
\end{figure}

\begin{proof}[Soluzione dell'esercizio]
\emph{Dominio e punto di accumulazione.}
La funzione
\[
  h(x):=x\sin\frac1x
\]
ha dominio \(A=\R\setminus\{0\}\). Il punto \(0\) è di
accumulazione per \(A\): per ogni \(\delta>0\), il
numero \(a=\delta/2\) appartiene ad \(A\) e soddisfa
\[
  0<|a|=\frac\delta2<\delta.
\]

\emph{Applicazione del teorema.}
Per ogni \(x\in A\), poiché
\(\lvert\sin(1/x)\rvert\le1\), si ha
\[
  0\le\left|x\sin\frac1x\right|
  =|x|\left|\sin\frac1x\right|\le|x|.
\]
Per ogni \(\varepsilon>0\), scegliendo
\(\delta=\varepsilon\), si verifica
\(|x|<\varepsilon\) quando \(0<|x|<\delta\).
Pertanto
\[
  \lim_{x\to0}|x|=0,
  \qquad
  \lim_{x\to0}(-|x|)=0.
\]
Poiché anche la funzione costante nulla ha limite
\(0\), applicando il \cref{thm:lez12-tre-funzioni}
alla prima disuguaglianza si ottiene
\[
  \lim_{x\to0}\left|x\sin\frac1x\right|=0.
\]
Dalla medesima stima segue inoltre
\[
  -|x|\le x\sin\frac1x\le|x|
  \qquad\forall x\in A.
\]
Applicando nuovamente il \cref{thm:lez12-tre-funzioni},
si conclude
\[
  \lim_{x\to0}x\sin\frac1x=0.
\]
\end{proof}

